\documentclass[10pt,a4paper]{amsart}
\usepackage[margin=19mm]{geometry}
\usepackage{amsmath,amssymb,mathtools}
\usepackage[scr=rsfs,cal=euler]{mathalfa}
\usepackage{xfrac}
\usepackage[unicode,hidelinks,bookmarks=false]{hyperref}
\newcommand{\ie}{i.e.\ }
\newcommand{\cf}{cf.\ }
\newcommand{\resp}{respectively\ }
\newcommand{\Subsection}{\subsection}
\providecommand{\nobreakdash}{\nobreak\hbox{-}}
\setlength{\emergencystretch}{2em}
\multlinegap0pt
\usepackage{bm}
\usepackage{bbm}
\usepackage{dsfont}
\usepackage{xspace}
\usepackage{cancel}
\usepackage{mathtools}
\usepackage{amsthm}
\makeatletter
\@ifundefined{thm}{\newtheorem{thm}{Thm}}{}
\@ifundefined{defn}{\newtheorem{defn}[thm]{Definition}}{}
\@ifundefined{lem}{\newtheorem{lem}[thm]{Lemma}}{}
\@ifundefined{prop}{\newtheorem{prop}[thm]{Proposition}}{}
\makeatother
\def\Kb{{\mathbf K}}
\def\NN{{\mathbb N}}
\def\RR{{\mathbb R}}
\def\VV{{\mathbb V}}
\def\d{\mathrm{d}}
\def\f{\frac}
\def\g{{\gamma}}
\def\sK{{\mathsf K}}
\def\sm{{\mathsf m}}
\def\sw{{\mathsf w}}
\newcommand{\wt}{\widetilde}
\title[Best approximation by polynomials on the conic domains]{Best approximation by polynomials on the conic domains}
\author{Yan Ge, Yuan Xu}
\date{Source: arXiv:2506.22916v1 (CC BY 4.0)}
\begin{document}
\maketitle

\noindent\emph{Source and licence.} Best approximation by polynomials on the conic domains. Version arXiv:2506.22916v1 (CC BY 4.0), distributed under the Creative Commons Attribution 4.0 International licence, as recorded in the arXiv licence metadata for this exact version and shown on the abstract page https://arxiv.org/abs/2506.22916v1. Full rights notice: licenses/arxiv-2506.22916-CC-BY-4.0.txt.\par\smallskip

\noindent\textit{Editorial scope.} Counted unit: the Prop whose source label is \texttt{prop:K-equ} in the article (source0.tex lines 1917--1923, source label \texttt{prop:K-equ}), delivered together with the article's own complete original proof of it (source0.tex lines 1925--1953, one \texttt{proof} environment of 29 lines). No mathematics is written, repaired or re-derived: statement and proof are the article's own text line for line. Required context retained uncounted: eq:cs-K-func (lines 960--968); eq:XY2 (lines 1781--1784); prop:IntV0V (lines 1789--1796); def:KfuncV (lines 1894--1908). Every definition, display and label of those lines is retained unchanged. Omitted from this package and not redistributed: 1--959, 969--1780, 1785--1788, 1797--1893, 1909--1916, 1924--1924, 1954--2110. Nothing inside a retained excerpt is omitted or altered: every retained body is delivered exactly as the article's lines are written, so no omission receipt of any kind accompanies this package. Citations: the retained excerpts carry no citation command at all, so no bibliography entry of the article is transcribed and no reference is redistributed. Reference adaptation: none was needed and none was made; every reference inside the delivered unit resolves inside it, so no unresolved reference or citation is produced. Printed slips kept literal: no printed word, symbol, hypothesis or number of the counted statement or of its proof was altered. Layout and class: the article's own class is \texttt{amsart} with packages including amsmath, amsthm, amsfonts, amssymb, accents, enumerate, color, graphicx, hyperref, geometry; the wrapper is the lane's pinned \texttt{amsart} editorial wrapper at 10pt on A4, which restates the theorem-like environments the retained text needs and adds the article's own macro definitions that the retained text needs (\texttt{\textbackslash Kb}, \texttt{\textbackslash NN}, \texttt{\textbackslash RR}, \texttt{\textbackslash VV}, \texttt{\textbackslash d}, \texttt{\textbackslash f}, \texttt{\textbackslash g}, \texttt{\textbackslash sK}, \texttt{\textbackslash sm}, \texttt{\textbackslash sw}, \texttt{\textbackslash wt}), with extra wrapper packages (bm, bbm, dsfont, xspace, cancel, mathtools). The delivered numbering is the unit's own. Assets: no retained span contains a figure environment, a graphics-inclusion command, a PDF-page inclusion, a table or a TikZ picture; no image file is copied into this package, the package declares no asset dependency and no image file is redistributed with it. Licence: version arXiv:2506.22916v1 of this article is distributed under the Creative Commons Attribution 4.0 International licence, as recorded in the arXiv licence metadata for this exact version and shown on the abstract page https://arxiv.org/abs/2506.22916v1. A full rights notice is delivered at licenses/arxiv-2506.22916-CC-BY-4.0.txt. Characters: every retained excerpt is pure ASCII, so the delivered engine, XeLaTeX with its default Latin Modern text font, reports no missing character.
\begin{defn}\label{eq:cs-K-func}
For $f\in L^p(\VV_0^{d+1}, \sw_{-1,\g})$, $1\leq p \le \infty$, and $r > 0$, the $K$-functional is defined by 
\begin{align*}
 \sK_r(f, h)_{p,\sw_{-1,\g}}: = \inf 
  	  \bigg \{ \|f-g\|_{p, \g}  & + h^r \left\|\varphi^r \partial^r_t g \right\|_{p, \g} 
	      + h^r  \max_{1\le i < j \le d} \left\| \f{1}{\sqrt{t}^r}D_{i,j}^r g\right\|_{p, \g} \bigg \},
\end{align*}
where the infimum is taken over all $g \in C^r(\VV_0^{d+1})$.
 \end{defn}  

\begin{equation}\label{eq:XY2}
	\tilde{f}\big(X, t\big)=\tilde{f}\big((x,x_{d+1}), t\big):=f(x,t), \quad  (x,t) \in \VV^{d+1}, \quad
	\big(X, t\big)\in\VV_0^{d+2},
\end{equation} 

\begin{lem}\label{prop:IntV0V}
Let $f: \RR^{d+1} \mapsto \RR$ be a continuous function. Let $x_{d+1} = \sqrt{t^2-\|x\|^2}$ for $\|x\| \le t$.
Then, with $X= (x,x_{d+1})$ and $X_* = (x,-x_{d+1})$, 
$$
	\int_{\VV_0^{d+2}} f(y,s) \sw_{-1,\g}(s) \d \sm(y,s) = \int_{\VV^{d+1}} 
	\big [ f\big(X,t\big) +  f\big(X_* ,t\big)\big ] W_{\g}(x,t) \d x\d t.
$$
\end{lem}

\begin{defn}\label{def:KfuncV} 
Let $f\in L^p(\VV^{d+1}, W_{\g,\mu})$, $1\leq p <  \infty,$ and $f\in C(\VV^{d+1})$
if $p=\infty$. Setting $\tilde{g}$ as the extension of $g$ defined in \eqref{eq:XY2}.
For   $r =1, 2$ and $\rho\geq 0$, the K-functional is defined by 
\begin{align*}
 & \Kb_r(f,\rho)_{p,W_{\g}} :=   \inf_{g \in C^r(\VV^{d+1})} 
	\bigg\{  \|f-g\|_{p,W_{\g}}  + \rho^r\left\|\varphi^{ r }\partial_t^rg  \right\|_{p,W_\g} \\ 
	& \qquad\qquad+ \rho^r \max_{1\leq i<j\leq d}   \Big \|\frac 1{\sqrt{t}^r}D_{i,j}^r g \Big \|_{p,W_{\g}}   
%	&\qquad\qquad 
	+ \rho^r \max_{1\leq i \leq d} \Big \|\frac 1{\sqrt{t}^r} \big ( \Phi \partial_i \big)^r g \Big \|_{L^p(\VV^{d+1},W_\g)}  
	%+ \rho^r\left\|\varphi^{ r }\partial_t^rg  \right\|_{p,W_\g} 
\bigg\},
\end{align*}
where $\varphi(x,t) = \sqrt{t(1-t)}$ and $\Phi(x,t) = \sqrt{t^2 - \|x\|^2}$. 
\end{defn} 

\begin{prop}\label{prop:K-equ}
Assume $f \in L^p(\VV^{d+1}, W_\g)$ if $1\le p < \infty$ and $f \in C(\VV^{d+1})$ if $p = \infty$. Let $\wt f$ be defined
as in \eqref{eq:XY2}. Then 
$$
 \Kb_r(f, \rho)_{L^p(\VV^{d+1}, W\g)} \sim  \sK_r \big(\wt f; \rho\big)_{L^p(\VV_0^{d+2}, \sw_{-1,\g})}.
$$
\end{prop}

\begin{proof}
Using Lemma \ref{prop:IntV0V}, we see that $\|f-g\|_{L^p(\VV^{d+1},W_\g)} = \|\wt f - \wt g\|_{L^p(\VV_0^{d+1}, \sw_{-1,\g})}$.  
Hence, comparing Definitions \ref{def:KfuncV} and \ref{eq:cs-K-func}, we only need to consider the term 
$\Big \|\frac 1{\sqrt{t}^r} \big ( \Phi \partial_i \big)^r g \Big \|_{L^p(\VV^{d+1},W_\g)}$ for $1 \le i \le d$ 
in $\Kb_r(f,\rho)_{p,W_{\g}}$. Let $(x,t) \in \VV^{d+1}$. Setting $X=(x,x_{d+1})$. Then $(X,t) \in \VV_0^{d+2}$ 
if $x_{d+1} = \Phi(x,t)$. For $r \in \NN$ and $(X,t) \in \VV_0^{d+2}$, we first claim that 
\begin{equation} \label{eq:FrXt}
F_r(X,t):=  (- \Phi \partial_i)^r g (x, t) = D_{i, d+1}^r \wt g(X ,t). 
\end{equation}
Since $D_{i,d+1}$ is an angular derivative, we obtain 
$$
   D_{i, d+1} \wt g(X,t) = (x_i \partial_{d+1} - x_{d+1} \partial_i) g(x,t) =   - x_{d+1} (\partial_i g)(y,t) = - \Phi(x,t) \partial_i g(x,t),
$$
which verifies \eqref{eq:FrXt} for $r =1$. Assume the identity has been proved for some $r\in \NN$. Then, taking the
derivative for $x_i$ and using $x_{d+1} = \Phi(x,t)$, we obtain 
\begin{align*}
 F_{r+1}(X,t)\,& =  - \Phi(x,t) \frac{\partial}{\partial x_i} F_r(X,t) \\
      & = - \Phi(x,t) \partial_i F_r(X,t) - \Phi(x,t) \partial_{d+1} F_r(X,t) \frac{\partial}{\partial x_i} \Phi(x,t) \\
      & = - x_{d+1} \partial_i F_r(X,t) + x_i \partial_{d+1} F_r(X,t) \\
      &  = D_{i, d+1} F_r(X,t) = D_{i, d+1}^{r+1} \wt g(X,t)
\end{align*}
by induction hypothesis, which completes the proof of \eqref{eq:FrXt}. With this identity, it follows readily, by 
Lemma \ref{prop:IntV0V}, that 
$$
  \Big \|\frac 1{\sqrt{t}^r} \big ( \Phi \partial_i \big)^r g \Big \|_{L^p(\VV^{d+1},W_\g)}
        =  \Big \|\frac 1{\sqrt{t}^r} D_{i,d+1}^r \tilde{g}\Big \|_{L^p(\VV_0^{d+2},\sw_{-1,\g})}.
$$
This proves the equivalence of the two $K$-functionals. 
\end{proof}

\end{document}
